\documentclass[
  12pt,
  oneside,
  openright,
  a4paper,
  english,
  brazil
]{abntex2}

\usepackage[T1]{fontenc}
\usepackage[utf8]{inputenc}
\usepackage{lmodern}
\usepackage{indentfirst}
\usepackage{microtype}
\usepackage{graphicx}
\usepackage{booktabs}
\usepackage{longtable}
\usepackage{array}
\usepackage{amsmath}
\usepackage{amssymb}
\usepackage{siunitx}
\usepackage{float}
\usepackage{hyperref}
\usepackage[brazilian,hyperpageref]{backref}
\usepackage[alf]{abntex2cite}

\title{Desenvolvimento de uma BMS 4S com ESP32-S3 para monitoramento, proteção e balanceamento de baterias de lítio}
\author{Seu Nome Completo}
\local{Sua Cidade}
\date{2026}

\instituicao{%
  Nome da Universidade\\
  Nome do Centro ou Faculdade\\
  Nome do Curso
}

\tipotrabalho{Trabalho de Conclusão de Curso}
\preambulo{Monografia apresentada ao curso de X da Y como requisito parcial para obtenção do título de Z.}

\orientador{Prof. Nome do Orientador}
\coorientador{Prof. Nome do Coorientador}

\hypersetup{
  pdftitle={\thetitle},
  pdfauthor={\theauthor},
  pdfsubject={BMS 4S com ESP32-S3},
  pdfcreator={LaTeX com abnTeX2},
  colorlinks=true,
  linkcolor=black,
  citecolor=black,
  urlcolor=blue
}

\makeindex

\begin{document}
\selectlanguage{brazil}
\frenchspacing

\pretextual

\imprimircapa
\imprimirfolhaderosto*

% Ficha catalografica:
% Substitua este bloco pela ficha emitida pela biblioteca da sua instituicao, se exigido.
% \begin{fichacatalografica}
%   \vspace*{\fill}
%   \hrule
%   \begin{center}
%     Ficha catalografica fornecida pela biblioteca
%   \end{center}
%   \hrule
% \end{fichacatalografica}

% Folha de aprovacao:
% Recomenda-se adaptar ao modelo institucional.
% \begin{folhadeaprovacao}
%   \begin{center}
%     \ABNTEXchapterfont\large\imprimirautor
%   \end{center}
%   \vspace*{\fill}\vspace*{\fill}
%   \begin{center}
%     \ABNTEXchapterfont\bfseries\large\imprimirtitulo
%   \end{center}
%   \vspace*{\fill}
%   \hspace{.45\textwidth}
%   \begin{minipage}{.5\textwidth}
%     \imprimirpreambulo
%   \end{minipage}
%   \vspace*{\fill}
%   Trabalho aprovado. \imprimirlocal, \imprimirdata:
%   \assinatura{\textbf{\imprimirorientador} \\ Orientador}
% \end{folhadeaprovacao}

\begin{resumo}
Este trabalho apresenta a proposta de desenvolvimento de um sistema de gerenciamento de baterias
para um pack de lítio em configuração \textit{4S}, utilizando um microcontrolador ESP32-S3 como
elemento de supervisão, telemetria e controle de alto nível. A pesquisa aborda os requisitos de
proteção do pack, a necessidade do uso de um \textit{Analog Front End} para medição segura das
células, a definição da máquina de estados da BMS e a organização do firmware para tratamento de
falhas, balanceamento passivo e configuração do sistema. A metodologia combina revisão bibliográfica,
análise de arquiteturas de mercado, modelagem da solução e estruturação de uma base inicial de
firmware e documentação técnica. Como resultado, propõe-se uma arquitetura coerente para uma BMS
4S com foco em proteção, expansibilidade e futura implementação em placa dedicada.

\textbf{Palavras-chave}: BMS. ESP32-S3. AFE. Bateria de lítio. Balanceamento. Proteção.
\end{resumo}

\begin{resumo}[Abstract]
\begin{otherlanguage*}{english}
This work presents the development proposal of a battery management system for a \textit{4S}
lithium battery pack, using an ESP32-S3 microcontroller as the high-level supervision, telemetry,
and control element. The study addresses pack protection requirements, the need for an
\textit{Analog Front End} for safe cell measurement, the definition of the BMS state machine,
and the firmware organization for fault handling, passive balancing, and system configuration.
The methodology combines bibliographic review, analysis of market architectures, solution modeling,
and the structuring of an initial firmware and technical documentation base. As a result, a
coherent architecture for a 4S BMS is proposed, focused on protection, expandability, and future
implementation on a dedicated printed circuit board.

\textbf{Keywords}: BMS. ESP32-S3. AFE. Lithium battery. Balancing. Protection.
\end{otherlanguage*}
\end{resumo}

\pdfbookmark[0]{\contentsname}{toc}
\tableofcontents*
\cleardoublepage

\textual

\chapter{Introdução}

\section{Contextualização}

O uso de baterias de íons de lítio em sistemas embarcados, mobilidade elétrica, armazenamento
de energia e equipamentos portáteis tem crescido de forma significativa nas últimas décadas.
Com esse crescimento, aumenta também a necessidade de sistemas de gerenciamento de baterias
(\textit{Battery Management Systems} -- BMS) capazes de monitorar, proteger e operar os packs
de forma segura \cite{andrea2010,plett2015v1}.

\section{Problema de pesquisa}

Em packs multicélula, como no caso de uma bateria \textit{4S}, a simples leitura da tensão total
do pack não é suficiente para garantir segurança operacional. É necessário monitorar cada célula
individualmente, acompanhar temperatura, corrente e estados de falha, além de aplicar estratégias
de proteção e balanceamento.

\section{Objetivo geral}

Desenvolver a proposta de uma BMS para bateria de lítio em configuração \textit{4S}, utilizando
um ESP32-S3 como unidade de supervisão e um \textit{Analog Front End} dedicado para medição e
proteção primária do pack.

\section{Objetivos específicos}

\begin{itemize}
  \item identificar os requisitos de proteção de uma BMS para pack 4S;
  \item analisar o papel do AFE em sistemas de gerenciamento de baterias;
  \item propor a arquitetura de hardware da solução;
  \item estruturar a máquina de estados e a lógica de proteção do firmware;
  \item organizar uma base de documentação e configuração para futura implementação física.
\end{itemize}

\section{Justificativa}

A relevância do tema se apoia na necessidade de segurança, robustez e rastreabilidade em sistemas
com células de lítio. Uma BMS mal projetada pode causar degradação precoce do pack, perda de
desempenho e, em casos severos, risco elétrico e térmico.

\section{Metodologia}

Este trabalho adota uma abordagem de pesquisa aplicada, de natureza tecnológica, combinando:

\begin{itemize}
  \item revisão bibliográfica sobre BMS, AFE e proteção de packs de lítio;
  \item análise de arquiteturas e componentes de mercado;
  \item modelagem da solução para uma BMS 4S;
  \item estruturação inicial do firmware e da documentação técnica.
\end{itemize}

\section{Estrutura do trabalho}

Este trabalho está organizado em sete capítulos. Após esta introdução, o Capítulo 2 apresenta
a fundamentação teórica. O Capítulo 3 descreve os requisitos e a metodologia. O Capítulo 4
apresenta a arquitetura proposta. O Capítulo 5 detalha o desenvolvimento da solução. O
Capítulo 6 discute os resultados e limitações. Por fim, o Capítulo 7 apresenta a conclusão.

\chapter{Fundamentação teórica}

\section{Baterias de lítio em configuração 4S}

Uma bateria em configuração \textit{4S} possui quatro células em série. Nesse arranjo, a tensão
total do pack resulta da soma das tensões individuais das células, o que exige monitoramento
por célula para evitar condições de sobretensão e subtensão.

\section{Battery Management System}

Um BMS é responsável por monitorar variáveis elétricas e térmicas do pack, proteger a bateria,
gerenciar carga e descarga e manter a operação dentro de limites seguros \cite{andrea2010}.

\section{Proteção em BMS}

As proteções mais relevantes em BMS incluem:

\begin{itemize}
  \item sobretensão por célula;
  \item subtensão por célula;
  \item sobrecorrente de carga;
  \item sobrecorrente de descarga;
  \item curto-circuito;
  \item sobretemperatura e subtemperatura.
\end{itemize}

Segundo Andrea \cite{andrea2010}, a robustez da proteção depende de uma combinação de hardware
dedicado e lógica de controle, especialmente em sistemas multicélula.

\section{Balanceamento de células}

O balanceamento busca reduzir a divergência entre tensões individuais das células. Entre as
estratégias possíveis, o balanceamento passivo é uma das mais comuns em aplicações de menor
complexidade, por dissipar energia das células mais carregadas em resistores de bleed
\cite{plett2015v2}.

\section{Analog Front End em BMS}

O \textit{Analog Front End} é o bloco responsável pela leitura segura de tensões, correntes e
temperaturas, além de frequentemente incorporar mecanismos de detecção de falhas e suporte a
balanceamento. Em sistemas reais, o AFE reduz a dependência do ADC do microcontrolador e
melhora a confiabilidade da proteção \cite{ti_bq76920,ad_passive_balancing}.

\section{Microcontroladores na supervisão da BMS}

O microcontrolador pode assumir as tarefas de:

\begin{itemize}
  \item supervisão de alto nível;
  \item telemetria;
  \item máquina de estados;
  \item registro de falhas;
  \item comunicação com interfaces externas.
\end{itemize}

No contexto deste trabalho, o ESP32-S3 foi escolhido por sua capacidade de processamento,
interfaces de comunicação e potencial de conectividade sem fio.

\chapter{Requisitos e metodologia do projeto}

\section{Requisitos funcionais}

Para a proposta da BMS 4S, consideram-se como requisitos funcionais:

\begin{itemize}
  \item leitura da tensão individual das quatro células;
  \item leitura de corrente do pack;
  \item leitura de temperatura por sensores NTC;
  \item controle independente dos caminhos de carga e descarga;
  \item suporte a balanceamento passivo;
  \item disponibilização de telemetria do pack.
\end{itemize}

\section{Requisitos de segurança}

Como requisitos de segurança, destacam-se:

\begin{itemize}
  \item bloqueio por sobretensão;
  \item bloqueio por subtensão;
  \item bloqueio por sobrecorrente;
  \item resposta segura em condição de curto-circuito;
  \item bloqueio por sobretemperatura;
  \item comportamento \textit{fail-safe} em caso de falha crítica.
\end{itemize}

\section{Parâmetros de configuração}

Os parâmetros iniciais do sistema são organizados em uma estrutura de configuração contendo:

\begin{itemize}
  \item características do pack;
  \item parâmetros de medição;
  \item limites de proteção;
  \item estratégia de balanceamento;
  \item convenções de telemetria.
\end{itemize}

\section{Metodologia de desenvolvimento}

A metodologia adotada foi composta por:

\begin{enumerate}
  \item levantamento bibliográfico;
  \item análise de AFEs e arquiteturas de mercado;
  \item definição dos requisitos do sistema;
  \item modelagem da arquitetura de hardware;
  \item modelagem da máquina de estados;
  \item implementação da base inicial de firmware.
\end{enumerate}

\chapter{Arquitetura da BMS 4S}

\section{Visão geral do sistema}

A arquitetura proposta divide a solução em três camadas principais:

\begin{itemize}
  \item camada de medição e proteção primária;
  \item camada de controle e supervisão;
  \item camada de potência e comutação.
\end{itemize}

\section{Arquitetura de hardware}

Os blocos de hardware propostos são:

\begin{itemize}
  \item pack de células 4S;
  \item AFE dedicado;
  \item resistor shunt para medição de corrente;
  \item sensores NTC;
  \item MOSFET de carga;
  \item MOSFET de descarga;
  \item resistores de balanceamento;
  \item fusível.
\end{itemize}

\section{Papel do AFE}

O AFE ocupa posição central na arquitetura, pois fornece medição multicélula, suporte à proteção
e interface segura com o microcontrolador. Para uma arquitetura 4S, famílias como a BQ76920
apresentam bom alinhamento com os requisitos de medição e proteção \cite{ti_bq76920}.

\section{Papel do ESP32-S3}

O ESP32-S3 foi definido como controlador de alto nível, responsável por:

\begin{itemize}
  \item ler a telemetria consolidada;
  \item executar a máquina de estados;
  \item tratar falhas;
  \item controlar os caminhos de potência;
  \item expor interfaces de monitoramento e configuração.
\end{itemize}

\section{Estágio de potência}

O estágio de potência prevê caminhos distintos para carga e descarga. Essa separação facilita
o bloqueio seletivo de eventos de falha conforme o tipo de condição detectada.

\section{Sensoriamento de corrente e temperatura}

O sensoriamento de corrente é realizado por resistor shunt, enquanto a monitoração térmica
utiliza, no mínimo, dois termistores NTC: um próximo às células e outro próximo ao estágio
de potência.

\chapter{Desenvolvimento do sistema}

\section{Estrutura do firmware}

O firmware foi organizado com separação entre tipos, configuração, monitoramento de bateria e
lógica de controle. Essa divisão favorece manutenção, evolução e futura substituição do modo
\textit{mock} por drivers reais de hardware.

\section{Máquina de estados}

A máquina de estados foi definida com os seguintes estados:

\begin{itemize}
  \item \textit{Startup};
  \item \textit{Idle};
  \item \textit{Charging};
  \item \textit{Discharging};
  \item \textit{Balancing};
  \item \textit{Fault}.
\end{itemize}

Essa modelagem torna explícita a relação entre telemetria, falhas, saídas e transições do
sistema.

\section{Lógica de proteção}

As proteções tratadas na base inicial do firmware incluem:

\begin{itemize}
  \item sobretensão por célula;
  \item subtensão por célula;
  \item sobrecorrente de carga;
  \item sobrecorrente de descarga;
  \item bloqueio térmico para carga e descarga;
  \item falha de AFE.
\end{itemize}

\section{Balanceamento}

O projeto adota, em sua primeira versão, balanceamento passivo. A lógica considera limiar de
início de balanceamento e delta mínimo entre a célula de maior e a de menor tensão.

\section{Estrutura de configuração em JSON}

Os parâmetros do sistema foram organizados em arquivo JSON para permitir rastreabilidade e futura
ligação com configuração persistente do firmware. Essa estrutura inclui parâmetros do pack,
medição, proteção, balanceamento, saídas e telemetria.

\chapter{Resultados e discussão}

\section{Cenários avaliados}

Nesta etapa inicial, os resultados concentram-se na modelagem e na estruturação do sistema.
Os principais cenários considerados foram:

\begin{itemize}
  \item carga com pack saudável;
  \item descarga com pack saudável;
  \item condição de desbalanceamento;
  \item sobretensão por célula;
  \item subtensão por célula;
  \item condição térmica crítica.
\end{itemize}

\section{Resultados obtidos}

Como resultado do desenvolvimento, foram obtidos:

\begin{itemize}
  \item definição de arquitetura para uma BMS 4S;
  \item estrutura de configuração do sistema;
  \item documentação técnica da proteção e do AFE;
  \item modelagem da máquina de estados;
  \item base inicial de firmware para futura integração com hardware.
\end{itemize}

\section{Análise crítica}

Embora a base proposta seja consistente para a etapa de projeto, a solução ainda depende de:

\begin{itemize}
  \item escolha final do AFE;
  \item esquemático completo;
  \item implementação do driver real do hardware;
  \item validação experimental em bancada;
  \item validação de proteção rápida em condição de falha.
\end{itemize}

\chapter{Conclusão}

Este trabalho organizou a base conceitual, arquitetural e de firmware para o desenvolvimento de
uma BMS 4S com ESP32-S3. A análise mostrou que o uso de um AFE dedicado é essencial para medição
segura das células e para uma estratégia de proteção mais robusta.

A proposta estruturada integra:

\begin{itemize}
  \item requisitos de proteção;
  \item arquitetura de hardware;
  \item máquina de estados;
  \item modelo inicial de firmware;
  \item rastreabilidade de configuração e documentação.
\end{itemize}

Como trabalhos futuros, destacam-se:

\begin{itemize}
  \item escolha final do AFE;
  \item desenvolvimento do esquemático e da PCB;
  \item implementação do driver de comunicação com o AFE;
  \item testes em bancada com fonte limitada;
  \item expansão da telemetria e da validação experimental.
\end{itemize}


\postextual

\bibliography{referencias}

\begin{apendicesenv}
\partapendices
\chapter{Checklist para adaptação institucional}

Antes da entrega final, recomenda-se validar:

\begin{itemize}
  \item modelo de capa da universidade;
  \item exigência de folha de aprovação;
  \item ficha catalográfica;
  \item regras próprias para resumo, abstract e listas;
  \item padrão institucional para referências e citações.
\end{itemize}
\end{apendicesenv}

\begin{anexosenv}
\partanexos
\chapter{Estrutura atual da documentação do projeto}

Os documentos técnicos do projeto podem ser usados como base para a redação do TCC:

\begin{itemize}
  \item \texttt{docs/arquitetura-bms-4s.md}
  \item \texttt{docs/protecao-bms.md}
  \item \texttt{docs/maquina-de-estados-bms.md}
  \item \texttt{docs/afe-em-bms.md}
  \item \texttt{config/bms\_config.json}
\end{itemize}
\end{anexosenv}

\printindex

\end{document}
